%%%%%%%%%%%%%%%%%%%%%%%%%%%%%%%%%%%%%%%%%%%%%%%%%%%%%%%%%%%%%%%%%%%%%%%%%%%%%%%%
%%  CONCLUSION
%%%%%%%%%%%%%%%%%%%%%%%%%%%%%%%%%%%%%%%%%%%%%%%%%%%%%%%%%%%%%%%%%%%%%%%%%%%%%%%%
\chapter{Conclusion}\label{conclusion}
We set out to show that a locally hosted model, grounded in a structured set of basic operations, can translate natural-language instructions into reliable, coordinated behavior across two robot arms without specialized training or hard-coded sequencing. The main takeaway is that our results support that claim across a broad range of tasks, while also marking the limits of the current implementation.

% ------------------------------------------------------------------------------
\section{What the System Accomplishes}
The system’s central pipeline works end-to-end. A single instruction flows through \gls{rag}-augmented retrieval, \gls{vlm} plan generation with chain-of-thought grounding, sequence-executor dispatch with variable-passing resolution, and physical execution in Unity. With the \gls{magi} backend, single-robot tasks achieved a 1.000 success rate across all 25 runs, and the B6 handoff across all 5 runs, with zero hallucinated operations and zero retries. These results show that the system can reliably turn language into coordinated action, and the dominant cost is contact-rich physical manipulation.

The pipeline composes its coordination at the planning layer rather than in the executor. Signal/wait synchronization points, parallel-group assignments for independent sub-tasks, and re-detection before each manipulation step are all generated by the \gls{vlm}; none is scripted into the executor, though the default prompt supplies rules that guide them. B10 illustrates this clearly, as the \gls{vlm} identified task-level independence between two disjoint pick-and-place subtasks and scheduled the grasp, field-detection, and placement phases into three concurrent groups while serializing the initial detections.

B8 stress-tests durability over a sustained multi-cycle sequence, and execution was flawless across all five runs with no cascading errors. The one run that scored as a failure missed only the benchmark’s exact chain-match check, despite identical execution. The main takeaway is sustained execution reliability, while op-level correctness and exact-sequence fidelity are scored separately.

Concurrent physical coupling marks the current capability boundary. Every model fails the simultaneous lift. The exposed failures are mis-realized cross-robot synchronization barriers and partner-aware target selection at the lift. The co-grasp itself is additionally blocked by the simulator’s single-parent attachment model (Sections~\ref{sec:results-dual} and~\ref{sec:bimanual}). This is the clearest remaining limitation.

Among the ablations, \gls{ros}/MoveIt routing showed the largest task-success gap, followed by reflection and negotiation. Reflection generalizes, helping or leaving the backend neutral, while negotiation helps only the default one; B16's gap, by contrast, is model-independent, tracing to \texttt{grasp\_object} timeouts that occur only under direct Unity routing, never under \gls{ros} for any model. On the grasp axis, the \gls{vgn} grasp path shows the largest and most general reliability effect, holding across all backends; an execution-path audit attributes it to grasp positioning and not \gls{vgn}'s predicted orientation (Section~\ref{sec:results-vgn}), whose own contribution awaits the depth-sensing upgrade below. \gls{rag} is model-dependent and, on the current task set, nets negative more often than it helps, and the \gls{kg} carries a small, consistent cost.

% ------------------------------------------------------------------------------
\section{Key Contributions}
The five contributions previewed in Section~\ref{sec:novelty} are realized and substantiated by the evaluation. The retrieval-grounded operation pipeline, a locally hosted reasoning model over a hierarchical 20-operation registry with \gls{rag} retrieval, an RT-H-style~\cite{belkhale2024rthactionhierarchiesusing} two-phase decomposition, and runtime variable-passing, decouples planning from execution without fine-tuning and underpins every benchmark in the suite. On top of that, the model coordinates the dual arms at the planning layer. The three optional layers each contribute a measurable, separately evaluated effect: consensus-gated negotiation, operation-gated reflection, and the AutoRT~\cite{ahn2024autortembodiedfoundationmodels} adaptation that pairs a structured object manifest with a hard kinematic safety layer the semantic filter cannot override. Together, the benchmarks confirm the central claim. Coordinated dual-arm behavior can be composed at runtime using a shared vocabulary of operations.

% ------------------------------------------------------------------------------
\section{Path to Deployment and Future Work}
The hardware adapter boundary analyzed in Section~\ref{sec:sim-to-real} enables the switch to hardware. The \texttt{--env real} flag selects the real-robot adapters, and the \gls{ros}/MoveIt path, \gls{ros}-TCP bridge, topic namespaces, and camera capture are already configured for dual-robot operation. What remains is the integration and validation: executing the planned trajectories on the physical AR4 \gls{ros} driver, calibrating the real cameras, and closing two sensing gaps, namely depth sensing, on which \gls{vgn}-quality grasp prediction depends via the existing \texttt{get\_depth\_frame} hook, and force-torque feedback, which would replace the contact-duration heuristic with ground-truth contact measurement.

Beyond those two hardware gaps, several further directions remain, following from the benchmark findings and the limitations identified in the last chapter.

\paragraph{Depth sensing upgrade.} Replacing the stereo pair with an active-stereo or time-of-flight depth sensor would make \gls{vgn}-based grasp planning reliable by supplying the dense, millimeter-level per-pixel depth that the current stereo reconstruction degrades on. It would also remove the current reliance on the simulator's \texttt{WorldState} oracle for object position, dimensions, and orientation (Section~\ref{sec:sim-to-real}), which a physical installation cannot provide, making this the gating prerequisite for real-world grasp localization.

\paragraph{Persistent spatial mapping.} Perception is currently memoryless. Every detection rebuilds the scene from a single stereo pair, with no spatial model carried across viewpoints or over time. Fusing the depth source above into a persistent volumetric map using Simultaneous Localization and Mapping (SLAM) or \gls{tsdf} fusion would give the planner object permanence under occlusion, multi-view grasp reconstruction, and a consistent workspace frame shared by both arms. It is also a prerequisite for moving beyond a fixed base.

\paragraph{Partner-aware target selection.} When B7’s execution reaches the concurrent lift, it fails because each arm resolves its own end-effector target in isolation, so two individually valid targets can land inside the 0.2~m separation check. The direct fix is a compatibility check at the lift. Before either \texttt{move\_to\_coordinate} commits, the coordination layer validates the two simultaneous targets against each other and nudges one apart, or serializes the pair, when they would violate the separation constraint.

\paragraph{Parse-level input validation.} The pre-parse guard from Section~\ref{sec:parse-validation}, extending the Robot Constitution’s code-level checks with robot-identifier and object-existence tests, would catch B9-class invalid inputs before any model call, regardless of model capability. Applying it to externally submitted commands is the next step.

\paragraph{Knowledge graph context targeting.} The B14 result motivates a more targeted context-injection design. Rather than prepending a full reachability summary, the \gls{kg} should supply only the spatial fact relevant to the planned operation. A retrieval step over the \gls{kg} predicates, analogous to the operation-level \gls{rag} step, would map the current planning context to the one or two graph edges most likely to inform the upcoming decision.

\paragraph{Real AR4 validation.} A complete evaluation on the physical AR4 arms would quantify the sim-to-real gap across all three system tiers: perception quality, trajectory execution fidelity, and grasp reliability.

\paragraph{Faster local inference.} The parsing latency, a median of roughly 20~s per command with the 24B reasoning model, is the system’s most persistent responsiveness constraint. Smaller distilled models fine-tuned on the operation registry vocabulary could narrow this window while maintaining the semantic grounding that enables zero-shot generalization. Speculative parsing of the likely follow-on operation during physical execution would further amortize inference costs by overlapping planning with motion execution.

% ------------------------------------------------------------------------------
\section{Closing Remarks}
The experimental results support our central claim. An off-the-shelf model can coordinate autonomous cooperation between two robotic arms by breaking natural-language tasks into basic operations. Single-robot manipulation and sequential dual-robot coordination are solved; concurrent bimanual manipulation is not, and the gap is specific. The \gls{vlm} already expresses and schedules physical simultaneity for independent sub-tasks (B10); on a shared object, it must additionally implement cross-robot synchronization barriers and choose mutually compatible targets for two actions that run at once, and B7 shows it does neither reliably. Those missing capabilities, dependable barrier realization, and partner-aware target selection under a shared physical constraint, define the boundary between the cooperation demonstrated here and genuine co-manipulation.

We built a modular architecture with a 20-operation registry, a \gls{rag}-augmented parser, a reflection-capable executor, an optional negotiation layer, and a hardware adapter boundary, enabling extensions at each of these points. The framework described here provides a workable foundation for deploying \gls{vlm}-driven cooperative robotics beyond simulation.

